\chapter{Related Work}\label{C:related}

Testing is a critical part to any software engineering process, not only to stop incidents stemming from the the product, but also partially related to the increase in popularity of agile methodologies \cite{chaos}. Many of these methodologies employ test driven development and continuous integration resulting in testing becoming more important throughout the development process. For these reasons there have been research papers that examine the different approaches that can be taken to identify redundant test cases and reduce the size of test suites \cite{wong1995effect, wong1999test, rothermel1998empirical, rothermel2002empirical,koochakzadeh2009test,zhang2011empirical,li2008static}.

It is unclear whether programmatically reducing a test suite's size is worth the tradeoff in the ability to locate bugs.  Wong et al. \cite{wong1995effect, wong1999test} found that test suite reduction does not severely impact fault detection capability. In contrast, Rothermel et al. \cite{rothermel1998empirical, rothermel2002empirical} found that test-suite reduction can severely impact the fault detection capability. This is the reason that our research uses human intuition to have the final say.

A popular technique used in detecting redundancy involves analysing the statement executions, known as coverage information. Maurer, Garousi and Koochakzadeh \cite{koochakzadeh2009test} attempt to answer the question, is coverage based identification enough information to determine redundant test cases? To identify a redundant test case, they state it as being a test case that does not improve a specific criterion. For example, in Figure \ref{fig:venndiagram} we see that according to the coverage data of test cases, T1 ... T5, T4 and T5 are fully redundant as T3 covers the statements that are executed by those tests. This criterion is examining statement coverage. They looked at two other criterion, branch coverage and granularities. Granularity criterion involved splitting the tests into setup, exercise (execution), verify (assert) and lastly teardown then performing analysis over each section. They implemented two different metrics in which both used the criterion described above. These two metrics examined redundancy on at a 1-1 and 1-all relationships. By comparing with manual inspection, they were able to determine the level of false positive and actual redundant tests. Of the redundant tests identified manually, the metrics matched 95\% of those. However, of the tests that were manually identified as being non-redundant, 52\% of these were identified as being redundant by the metrics. They concluded that coverage-based information is vulnerable in giving false-positives when identifying redundant test cases, suggesting common code paths as being the main protagonist. 

\begin{figure}[h]
\begin{center}
\includegraphics[]{VennDiagram.png}
\end{center}
\caption{A Typical Criterion}
\label{fig:venndiagram}
\end{figure}

A technique that Zhang, Marinov, Zhang and Khurshid \cite{zhang2011empirical} examined was the use of a greedy technique in comparison to heuristics with the coverage information. This required a test requirement to be set by the tester, for example using statement coverage. The greedy technique would greedily select a test case that satisfies the maximum number of unsatisfied test requirements and would continue until all the test requirements had been satisfied. So the new test suite will contain exactly the same coverage as the old test suite while removing redundant test cases. This means that if a test subsumed another, then it would always be removed. The heuristic implementation was first conceived by Harrold, Gupta and Soffa \cite{harrold1993methodology} where essential test cases are selected as early as possible. Essential being that only when one test case satisfies a test requirement exclusively. The heuristic approach resulted in the most cost-effective reduction showing that although greedy approach worked, there were better techniques available.

In situations where it is not possible to generate a spectrum to analyse, static analysis can be used to determine the level of redundancy. Robinson, Li and Francis \cite{li2008static} examine the possibility to do this. For the benchmark they use, the test cases are written in a high level automation framework and consist of a list of commands. The commands perform actions such as file copying and loading configurations. To identify redundant test cases they examine the test cases commands as well as the instructions within the procedures that the test case loads. To calculate the similarities between two test cases, they consider three different metrics, Manhattan distance, unigram cosine similarity and bigram cosine similarity. They each were measuring how closely related two tests were based on the sequence of commands and procedures loaded. Their findings were similar to Maurer et al. \cite{koochakzadeh2009test} in that there are a large number of false - positives.

Static checking has several limitations in comparison to dynamic. Firstly, during run-time is when a large part of the paper trail is accessible. Before then the only information available are the method calls. Secondly, it is more difficult to analyse, due to inheritance in many languages adding a complex situation when determining method execution. Static checking therefore would provide useful information in a framework where dynamic data can not be collected and allows for the test cases to be examined in a static fashion such as the one used by Robinson, Li and Francis \cite{li2008static}. Dynamic allows for more data to be collected such as the parameters passed to a method as well as the ability to be executed on more suites. The papers examined looked at the coverage of a test suite at several different criterion levels but none looked at the spectra explicitly. This leaves a potentially useful approach to the problem that may help determine the level of redundancy within a test suite. 